\documentclass[11pt]{article}
\usepackage[utf8]{inputenc}
\usepackage[T1]{fontenc}
\usepackage{lmodern}
\usepackage[margin=1in]{geometry}
\usepackage{amsmath,amssymb,mathtools}
\usepackage{graphicx}
\usepackage{grffile}
\usepackage{float}
\usepackage{hyperref}
\usepackage{caption}
\usepackage{parskip}
\usepackage{enumitem}
\usepackage{xurl}
\setlength{\parskip}{0.65em}
\setlength{\parindent}{0pt}
\begin{document}
\title{Daily Research Digest --- 2026-04-25}
\author{}
\date{}
\maketitle
\textbf{Topics:} galaxy clusters, galaxies, gravitational lensing, dark matter.
\section*{Synthesis}
Here's a summary and categorization of the recent astrophysics papers based on their primary focus areas:
\subsection*{Dark Matter}
1. **Multi-wavelength study of EP250416a / GRB 250416C**
\begin{itemize}
\item Focus: Analysis of the gamma-ray burst with multi-wavelength observations, including X-rays and optical data.
\end{itemize}
2. **Radiation properties of a regular black hole embedded in a Dehnen-type dark matter halo with a thin accretion disk**
\begin{itemize}
\item Focus: Investigating radiation properties, shadows, and geodesic structures around a regular black hole within a dark matter halo.
\end{itemize}
3. **Fermion Condensate Inflation, Dynamical Waterfall Mechanism and Primordial Black Holes**
\begin{itemize}
\item Focus: Exploring fermion condensate inflation and its role in forming primordial black holes as seeds for dark matter.
\end{itemize}
4. **SPURS: Bursty Star Formation in an Extremely Luminous Weak Emission Line Galaxy at z=9.3**
\begin{itemize}
\item Although primarily focused on galaxies, this paper discusses star formation in the context of galaxy evolution at high redshifts which can have implications for dark matter.
\end{itemize}
\subsection*{Galaxies}
1. **Signatures of Very Massive Stars in the Epoch of Reionization**
\begin{itemize}
\item Focus: Observing strong wind lines and emission from very massive stars using JWST to understand star formation in early galaxies.
\end{itemize}
2. **Multi-wavelength study of EP250416a / GRB 250416C**
\begin{itemize}
\item Focus: Detailed analysis of the burst's prompt emission, afterglow evolution, and physical origin.
\end{itemize}
3. **SPURS: Bursty Star Formation in an Extremely Luminous Weak Emission Line Galaxy at z=9.3**
\begin{itemize}
\item Focus: Ultra-deep spectroscopy to study star formation histories and properties of a super-luminous galaxy.
\end{itemize}
4. **Radiation properties of a regular black hole embedded in a Dehnen-type dark matter halo with a thin accretion disk**
\begin{itemize}
\item Focus: Examining radiation from an accretion disk around a black hole within a dark matter halo.
\end{itemize}
5. **Signatures of Very Massive Stars in the Epoch of Reionization**
\begin{itemize}
\item Focus: Using JWST to study star formation and properties of very massive stars at high redshifts.
\end{itemize}
6. **SPURS: Bursty Star Formation in an Extremely Luminous Weak Emission Line Galaxy at z=9.3**
\begin{itemize}
\item Focus: Understanding the star formation history of a super-luminous galaxy with weak emission lines.
\end{itemize}
7. **Signatures of Very Massive Stars in the Epoch of Reionization**
\begin{itemize}
\item Focus: JWST spectroscopy to understand the role of very massive stars in early galaxies.
\end{itemize}
8. **Fermion Condensate Inflation, Dynamical Waterfall Mechanism and Primordial Black Holes**
\begin{itemize}
\item Although primarily focused on inflation models, it discusses implications for dark matter through primordial black holes.
\end{itemize}
\subsection*{Summary}
\begin{itemize}
\item **Dark Matter**: Four papers (1, 2, 3, and partially 4).
\item **Galaxies**: Eight papers (5, 6, 7, 8, 9, 10, 11, and partially 4).
\end{itemize}
The studies highlight the interplay between dark matter and galaxy formation, emphasizing the role of very massive stars and their impact on early galaxy evolution. The use of JWST for deep spectroscopy is a common theme in understanding star formation at high redshifts.
\newpage
\section*{Paper Summaries and Figures}
\subsection*{1. An RXTE Search for the Sterile Neutrino Decay in Galaxy Clusters}
\textbf{Focus:} galaxy clusters, galaxies, dark matter\\
\textbf{Authors:} Mark J. Henriksen\\
\textbf{Published:} 2026-04-23T16:11:53+00:00\\
\textbf{PDF:} \href{https://arxiv.org/pdf/2604.21817v1}{https://arxiv.org/pdf/2604.21817v1}\\
\textbf{Summary:}
We have searched for the 3.55 keV line from sterile neutrino decay using 3.1 megaseconds of RXTE cluster data. A 2.5$\sigma $ excess of emission over a thermal model is found over the energy span of the 3.55 keV line in the combined spectra of the eight clusters that individually have an excess. The residuals are added to increase the signal to noise ratio of the excess, which is then modeled with a Gaussian to simulate the instrumental spectral response. We find a significant correlation (r = 0.76) for a line centered at 3.6 keV with a model flux of 3.07 x 10$^{-5}$ ph cm$^{-2}$ s$^{-1}$. Mixing angle for detected clusters ranges from 0.35 to 6.2 x 10$^{-10}$. The decay rate inferred from the line flux is strongly correlated (r = 0.87) with cluster temperature, which is due to hotter, more massive clusters having a larger amount of dark matter. Approximately half of the decay line total flux comes from the Coma cluster. We fit the Coma cluster spectrum with two different three-component models. The first includes a Gaussian fixed at 3.55 keV to model soft emission. The second three-component model uses a second thermal component to model soft emission. The model fit parameters indicate that the second thermal component is modeling high-energy residuals rather than low ones, where the Gaussian is. Though our line fluxes exceed most reported detections and upper limits, they do not overproduce the dark matter. We conclude that some fraction of the marginally detected excess could be attributed to the decay line since low-temperature thermal emission and systematics fail to model it completely.
\subsubsection*{Selected figures}
\begin{figure}[H]
\centering
\includegraphics[width=\linewidth,height=0.42\textheight,keepaspectratio]{/home/kf/research-assistant/data/cache/figures/http-arxiv.org-abs-2604.21817v1/figure\_1.png}
\captionsetup{labelformat=empty}\caption{Figure 3. The ratio (data/model) in each channel covering the energy range of the 3.55 k
feature folded through the RXTE PCA energy resolution is shown for each cluster. Lines a
that bracket the systematic error for the PCA as ±∼0.5\%. The features can be compared
systematic error and the number of c}
\end{figure}
\newpage
\subsection*{2. Informative Priors on Primordial Non-Gaussianity Bias $b_\phi $ From Galaxy Formation}
\textbf{Focus:} galaxies\\
\textbf{Authors:} Anne Moore, Lucia A. Perez, Elisabeth Krause\\
\textbf{Published:} 2026-04-23T15:46:53+00:00\\
\textbf{PDF:} \href{https://arxiv.org/pdf/2604.21790v1}{https://arxiv.org/pdf/2604.21790v1}\\
\textbf{Summary:}
Constraining primordial non-Gaussianity via its scale-dependent imprint on galaxy clustering requires knowledge of the bias parameter $b_\phi $, which is exactly degenerate with $f^{\rm{loc}}_{\rm{NL}}$ at leading order. To break this degeneracy, current analyses adopt the relation $\left(b_\phi  = 2\delta _c\left(b_1 - 1\right)\right)$ based on the assumption of a universal mass function. This relation is known to break down for physically motivated galaxy selections, introducing systematic errors in the inferred $f^{\rm{loc}}_{\rm{NL}}$ that scale directly with the assumed $b_\phi $ prior. We present a framework to construct physically motivated, observation-conditioned priors on $b_\phi $ by marginalizing over galaxy formation uncertainties. We use the CAMELS-SAM simulation suite, augmented by separate Universe simulations, to measure galaxy formation observables, like the stellar mass function (SMF) and the stellar-to-halo mass relationship (SHMR), and $b_\phi $ across a range of galaxy formation parameters. From these measurements, we construct a distribution of $b_\phi $ conditioned on observations, and we select our galaxy sample to resemble the DESI Emission Line Galaxy (ELG) sample. Conditioning on the SMF or SHMR decreases $\sigma _{b_\phi }$ from $0.69$ to $0.08$ and $0.02$ respectively -- reductions of $88\%$ and $97\%$ -- with consistent results when conditioning on the observed data directly. Despite substantial shifts in the galaxy formation posteriors driven by known SC-SAM discrepancies at high halo masses, the resulting $b_\phi $ distributions remain mutually consistent across all observables. The SMF and SHMR are found to carry sufficient constraining power to reduce the galaxy formation uncertainty in $b_\phi $ relevant for $f^{\rm{loc}}_{\rm{NL}}$ inference with next-generation spectroscopic surveys
\newpage
\subsection*{3. Accelerating scaling solutions from dark matter particle creation}
\textbf{Focus:} dark matter\\
\textbf{Authors:} Sudip Halder, Jaume de Haro, Supriya Pan, Emmanuel N. Saridakis, Tapan Saha, Subenoy Chakraborty\\
\textbf{Published:} 2026-04-23T15:28:34+00:00\\
\textbf{PDF:} \href{https://arxiv.org/pdf/2604.21770v1}{https://arxiv.org/pdf/2604.21770v1}\\
\textbf{Summary:}
This article opens new window to obtain accelerating scaling attractors without any need of dark energy. We study cosmological dynamics in a two-fluid system where pressureless dark matter (DM) undergoes adiabatic particle creation and exchanges energy with a barotropic fluid. Considering six widely used interaction prescriptions, we formulate the corresponding autonomous systems in a compact phase space and perform a unified dynamical analysis. We find that accelerating scaling attractors, namely late-time states where both fluids coexist with fixed energy fractions, arise only when the interaction is controlled by the DM density and energy flows from DM to the second fluid. Such attractors appear in the global and local DM-based interactions, and in the global mixed case, but are entirely absent when the interaction depends on the second fluid or on local mixed terms, which instead drive the universe to a DM-dominated accelerating phase. These results clarify the unique conditions under which matter creation can mimic dark-energy-like behaviour without introducing a dark-energy component.
\subsubsection*{Selected figures}
\begin{figure}[H]
\centering
\includegraphics[width=\linewidth,height=0.42\textheight,keepaspectratio]{/home/kf/research-assistant/data/cache/figures/http-arxiv.org-abs-2604.21770v1/figure\_1.png}
\captionsetup{labelformat=empty}\caption{FIG. 7. Summary of accelerating scaling attractors identif
in this study. Each row corresponds to a particular ene
density (green, yellow, or pink), while each column represe
the type of interaction: global rate (dark shading) or lo
rate (light shading).}
\end{figure}
\newpage
\subsection*{4. Delving into the depths of NGC 3783 with XRISM: V. Broad-band modeling of ionized outflows}
\textbf{Focus:} galaxies\\
\textbf{Authors:} Keqin Zhao, Jelle S. Kaastra, Liyi Gu, Missagh Mehdipour, Megan E. Eckart, Keigo Fukumura, Matteo Guainazzi, Chen Li, Christos Panagiotou, Matilde Signorini\\
\textbf{Published:} 2026-04-23T14:16:59+00:00\\
\textbf{PDF:} \href{https://arxiv.org/pdf/2604.21710v1}{https://arxiv.org/pdf/2604.21710v1}\\
\textbf{Summary:}
The Seyfert 1 galaxy NGC 3783 hosts a multiphase warm absorber (WA) that has been extensively studied in the X-ray band. High-resolution spectra from 2000-2001 revealed a complex outflow with multiple ionization and velocity components. Two decades later, new XMM-Newton and XRISM observations allow us to investigate the long-term evolution of these outflows. We perform joint spectral modeling of the XMM-Newton/RGS and XRISM/Resolve time-averaged spectra using the pion photoionization code within SPEX. We derive the ionization parameter, column density, turbulent velocity, and outflow velocity for each absorption component, and investigate their thermal stability and Absorption Measure Distribution (AMD) to characterize the physical and dynamical properties of the WA in NGC 3783 in 2024. We compare these results with the 2000-2001 epoch to assess long-term variability, stability, and possible changes in the absorber population. We identify eight WA components spanning log $ξ=$ 1.08-3.38 and outflow velocities of 480-1230 km s$^{-1}$. The ranges of column densities and turbulent velocities remain broadly consistent with the WAs from 2000-2001, but the earlier data contained more low-ionization, high-velocity components. The total column density in 2024 is 1.5 times larger than in 2000-2001, requiring replenishment by fresh material. The dominant Unresolved Transition Array (UTA) absorber (Comp. B3) has increased its column density by a factor of three while maintaining a similar ionization parameter. The WAs in NGC 3783 have undergone significant structural and dynamical evolution over the past 24 years.
\subsubsection*{Selected figures}
\begin{figure}[H]
\centering
\includegraphics[width=\linewidth,height=0.42\textheight,keepaspectratio]{/home/kf/research-assistant/data/cache/figures/http-arxiv.org-abs-2604.21710v1/figure\_1.png}
\captionsetup{labelformat=empty}\caption{Fig. 3. Relations between the parameters of the eight outflow compo-
nents derived from the XRISM Resolve spectrum of NGC 3783. The
top panel displays the column density NH as functions of the ionization
parameter (log ξ). The middle and bottom panels show the turbulent ve-
locity ($\textbackslash{}sigma$v) and outflow vel}
\end{figure}
\newpage
\subsection*{5. Saturation Mechanisms in the Interacting Dark Sector}
\textbf{Focus:} dark matter\\
\textbf{Authors:} Andronikos Paliathanasis, Kevin J. Duffy\\
\textbf{Published:} 2026-04-23T13:36:31+00:00\\
\textbf{PDF:} \href{https://arxiv.org/pdf/2604.21671v1}{https://arxiv.org/pdf/2604.21671v1}\\
\textbf{Summary:}
We introduce a family of phenomenological cosmological models featuring an interacting dark sector modulated by a sparseness scale parameter, in order to describe the late-time accelerated expansion of the universe. The sparseness scale, inspired by well-established saturation mechanisms in ecology and biology, is introduced in the interaction as a half-saturation constant that bounds the energy exchange between dark matter and dark energy, controls the dynamical behaviour of the physical variables and can prevent the phantom crossing. We consider three nonlinear interacting models, where two of them recover the linear interacting scenarios when the sparsity parameter vanishes. We examine the phase-space of the cosmological field equations by using the Hubble normalization approach. We determine the stationary points and their stability properties in order to reconstruct the asymptotics behaviour of the field equations. Such an analysis allows us to demonstrate the effects of the sparseness scale on the background dynamics. We test the interacting models with observational data. Specifically, we employ Supernovae catalogues, cosmic chronometers, Baryon Acoustic Oscillation measurements from DESI DR2, and redshift-space distortion measurements of the growth of large-scale structure through the $f$ and $f\sigma _8$ observables. The Bayesian analysis suggests that, for two of the three models, a vanishing sparsity parameter is disfavoured at more than the 95\textbackslash{}\% confidence interval, providing observational support for a nonzero sparseness scale in the dark sector interaction.
\newpage
\subsection*{6. Conformal anomaly transport induced by dark photon}
\textbf{Focus:} dark matter\\
\textbf{Authors:} Marek Rogatko, Karol I. Wysokinski\\
\textbf{Published:} 2026-04-23T13:27:00+00:00\\
\textbf{PDF:} \href{https://arxiv.org/pdf/2604.21664v1}{https://arxiv.org/pdf/2604.21664v1}\\
\textbf{Summary:}
We have considered the problem of the influence of inhomogeneity of gravitational field on transport effects predicted by the field theory describing massless Dirac fermions in the Maxwell and dark matter background. As a model of dark sector one takes into account dark photon model, where the hidden sector is described by the auxiliary U(1)-gauge field coupled to the visible sector. Elaborating the model we restrict our considerations to the case when Weyl type conformal transformation slightly differs from the Minkowski spacetime. This assumption simplifies the calculations and enables us not to use complicated methods of the quantum field theory in the curved background. The resulting currents stemming both from visible and dark sectors are proportional to the adequate beta functions appearing in the elaborated systems. For charge-less dark sector we predict corrections to the scale conductivities in both sectors: linear in $\textbackslash{}alpha$ in the dark sector and quadratic in the visible one.
\newpage
\subsection*{7. Constraining dark matter self-interaction from kinetic heating in neutron stars}
\textbf{Focus:} dark matter\\
\textbf{Authors:} Sambo Sarkar\\
\textbf{Published:} 2026-04-23T13:17:20+00:00\\
\textbf{PDF:} \href{https://arxiv.org/pdf/2604.21652v1}{https://arxiv.org/pdf/2604.21652v1}\\
\textbf{Summary:}
Dark matter search strategies have started advancing towards the neutrino fog. In this regard, compact objects such as neutron stars have already demonstrated their ability in probing such low DM-nucleon cross-sections from dark matter induced effects. In the optically thin limit, effect of dark matter self-interaction becomes relevant and may assist the capture and thermalization of dark matter inside stars, imparting observable changes on neutron star temperatures. The resulting radiation although weak can be potentially detected by the James Webb Space Telescope and upcoming Thirty Meter Telescope and the European Extremely Large Telescope. Observation of cold neutron stars accompanied by advancements in direct detection probes would provide stringent constraints or a smoking-gun signature for dark matter self-interactions. The potential detection of a neutron star with surface temperatures $\sim (1000 - 1200)$ K in the optically thin limit can push the bounds on asymmetric dark matter self-interaction cross-section to approximately two orders of magnitude more stringent than the bullet cluster.
\newpage
\subsection*{8. Aspects of gravitational clustering and structure formation in the Universe}
\textbf{Focus:} galaxies\\
\textbf{Authors:} Swati Gavas\\
\textbf{Published:} 2026-04-23T12:52:25+00:00\\
\textbf{PDF:} \href{https://arxiv.org/pdf/2604.21634v1}{https://arxiv.org/pdf/2604.21634v1}\\
\textbf{Summary:}
The distribution of galaxies, halo abundance, and peculiar velocities are influenced by non-linear gravitational interactions, making the study of non-linear evolution crucial for accurate cosmological predictions. We explore these aspects using N-body simulations. Theoretical models of the halo mass function (HMF) can be formulated without referencing a cosmological model or input power spectrum. HMF obtained from N-body simulations show systematic deviations of 5-20\textbackslash{}\% from theoretical predictions. The physical origin of deviations may result from cosmology, the power spectrum, or both. We examine HMF deviations from universality for scale-free power spectra with an Einstein-de Sitter cosmology. We demonstrate that the mass function exhibits an explicit dependence on the slope of the input power spectrum. We find that an effective index of the $Λ$CDM model can correspond to the HMF from scale-free cosmologies as a first approximation. Furthermore, structure formation has led to deviations from homogeneity and isotropy on scales up to at least $100$ Mpc/h, expected to affect measurements of $H_0$. We revisit this issue of the concordance model. We find a correlation between errors in $H_0$ estimates and the density around the observer. Further, our mock observations reveal that deviations of up to 5\textbackslash{}\% can occur in Milky Way-sized halos. While this finding alone does not fully resolve the Hubble tension, it may account for part of it. It is essential to understand the limitations of N-body simulations to avoid misinterpreting data. We show that the missing power at small scales introduces errors in the root-mean-square fluctuations and in the simulated mass function. Our analytical calculation indicates that mode coupling between small and large scales depends on resolving collapsed halos. Therefore, accurate mode coupling estimates require sufficient halos in the simulation.
\newpage
\subsection*{9. Multi-wavelength study of EP250416a / GRB 250416C: An Optically Dark Long GRB with a Late Jet Break}
\textbf{Focus:} galaxies\\
\textbf{Authors:} Guoying Zhao, Duo-Le Cao, Rong-Feng Shen, Hui Sun, Chi-Chuan Jin, Wei-Min Yuan, Chen-Wei Wang, Shao-Lin Xiong, Dmitry Svinkin, Dong Xu, Shuai-Qing Jiang, Peter G. Jonker, Yun Wang, Hao Zhou, Chang Zhou, Xinlei Chen, Kaushik Chatterjee, Xue-Feng Wu, Xiao-Feng Wang, Chun Chen, Yuan Liu, Andrew J. Levan, Jennifer Alexandra Chacon Chavez, Jonathan Quirola-Vásquez, Franz E. Bauer, Antonio Martin-Carrillo, Gregory Corcoran, Daniele B. Malesani, Dmitry Frederiks, Anna Ridnaia, Alexandra L. Lysenko, Mikhail Ulanov\\
\textbf{Published:} 2026-04-23T12:44:15+00:00\\
\textbf{PDF:} \href{https://arxiv.org/pdf/2604.21624v1}{https://arxiv.org/pdf/2604.21624v1}\\
\textbf{Summary:}
We present multi-wavelength study of the $\gamma $/X-ray transient EP250416a (also designated GRB 250416C), triggered by the Einstein Probe (EP) Wide-field X-ray Telescope and also by SVOM and Konus-Wind. Observations spanning the gamma-ray, X-ray, and optical bands facilitated detailed analysis of the burst's prompt emission, afterglow evolution, and physical origin. EP250416a exhibits a burst duration of 30 s in X-ray and 17.7 s in gamma-rays, with joint spectral fitting of 0.5-5000 keV data gives $E\rm_{peak}=342_{-232}^{+90}$ keV. Optical spectroscopy of the afterglow, acquired with the Gemini Multi-Object Spectrograph (GMOS) on Gemini South, yielded a redshift of $z=0.963$. Accounting for the measured redshift, the isotropic energies are $E\rm_{X,iso}=2.7_{-0.5}^{+0.9}\times10^{50}$ erg and $E\rm_{\gamma ,iso}=7.34_{-2.1}^{+5.1}\times10^{51}$ erg, aligning with the Amati relation for long GRBs. The fluence ratio $\rm S(25-50~keV)/S(50-100~keV)=0.78_{-0.15}^{+0.1}$ classifies EP250416a as an X-ray rich (XRR) GRB. The X-ray afterglow shows an initial shallow decay ($\alpha \approx -0.5$) transitioning to a canonical decay phase ($\alpha \approx -1$), with a very late jet break at $t\sim 1.5\times 10^6$ s, corresponding to a jet half-opening angle of $θ_j=10.6_{-1.8}^{+1.9}$ degrees. EP250416a is optically dark, as it shows only a faint $r$-band detection ($r=24.16$ mag) from Gemini South-GMOS and a low optical-to-X-ray spectral index $\beta _{\rm OX} = 0.3$. This may be attributed to significant host-galaxy extinction, with a required $A_V^{\text{host}}=5.5\ \text{mag}$ derived from the extinction curve model.
\subsubsection*{Selected figures}
\begin{figure}[H]
\centering
\includegraphics[width=\linewidth,height=0.42\textheight,keepaspectratio]{/home/kf/research-assistant/data/cache/figures/http-arxiv.org-abs-2604.21624v1/figure\_1.png}
\captionsetup{labelformat=empty}\caption{mage of the EP250416a captured by the EP/WXT CMOS1
chip at 17:53:59 UTC on 16 April 2025. the source is localize
256.42◦and Dec. = 25.78◦. tions of EP250416a : the first at ∼4.8 hours (PI: Kennea
scond at ∼20.5 hours (PI: EPSC) after the WXT trigge
bservations complemented the EP/FXT data by coverin}
\end{figure}
\newpage
\subsection*{10. Radiation properties of a regular black hole embedded in a Dehnen-type dark matter halo with a thin accretion disk}
\textbf{Focus:} dark matter\\
\textbf{Authors:} Tianyou Ren, Jing-Ya Zhao, Xiaomei Liu, Rong-Jia Yang\\
\textbf{Published:} 2026-04-23T12:39:49+00:00\\
\textbf{PDF:} \href{https://arxiv.org/pdf/2604.21615v1}{https://arxiv.org/pdf/2604.21615v1}\\
\textbf{Summary:}
We investigate the shadow, timelike geodesic structure, radiation properties of thin accretion disks, and optical appearance of a static spherically symmetric regular black hole, constructed based on the Dehnen-type density profile. Using observational data from M87* and Sgr A*, we constrain the model parameter $a$ at both $1\sigma $ and $2\sigma $ confidence levels. Based on the Page--Thorne model, we calculate the local radiative flux, redshift factor distribution, and the radiation flux received by a distant observer, systematically examining the effects of the parameter $a$ and the viewing angle on the black hole image. The results show that larger $a$ will enlarge the effective radiation area of the accretion disk and significantly enhance the asymmetry and Doppler boosting effects of the direct and secondary images at large viewing angles.
\newpage
\subsection*{11. Fermion Condensate Inflation, Dynamical Waterfall Mechanism and Primordial Black Holes}
\textbf{Focus:} dark matter\\
\textbf{Authors:} Stephon Alexander, Pisin Chen, Jinglong Liu, Antonino Marciano, Misao Sasaki, Xuan-Lin Su\\
\textbf{Published:} 2026-04-23T10:59:00+00:00\\
\textbf{PDF:} \href{https://arxiv.org/pdf/2604.21535v1}{https://arxiv.org/pdf/2604.21535v1}\\
\textbf{Summary:}
Fermion condensate inflation, where inflation emerges from four-fermion interactions induced by spacetime torsion, removes the need for additional scalar fields beyond the Standard Model. In this framework, the fermion field can be decomposed into two distinguished sectors, each giving rise to bound states. After integrating out fermions, the bound fields play the roles of the inflaton and the auxiliary fields, resembling hybrid inflation with a waterfall mechanism. The inclusion of an axial chemical potential naturally introduces a mechanism to end inflation and trigger instant preheating. During the waterfall phase, the effective potential of the fermion condensate supports the formation of non-topological solitons such as Q-balls, which act as seeds of primordial black holes. This model is intrinsically connected to Chern-Simons gravity, which implies a parity-violating universe. Consequently, both the primordial black hole (PBH) dark-matter abundance and parity-violation signatures could provide observational tests of the model.
\subsubsection*{Selected figures}
\begin{figure}[H]
\centering
\includegraphics[width=\linewidth,height=0.42\textheight,keepaspectratio]{/home/kf/research-assistant/data/cache/figures/http-arxiv.org-abs-2604.21535v1/figure\_1.png}
\captionsetup{labelformat=empty}\caption{FIG. 1: We plot the normalized effective potential
V (A, B)/Λ4, with the B-axis truncated at 0.15Λ for a
clearer view of the minima. The minima of A are
highlighted in red, and those of B in blue, the other
parameters being kept constant. Using the parameter
specified in Sec. (IV A), the figure show}
\end{figure}
\newpage
\subsection*{12. SPURS: Bursty Star Formation in an Extremely Luminous Weak Emission Line Galaxy at $z=9.3$}
\textbf{Focus:} galaxies\\
\textbf{Authors:} Zuyi Chen, Daniel P. Stark, Charlotte A. Mason, Adele Plat, Viola Gelli, Peter Senchyna, Keerthi Vasan G. C., Ryan Endsley, Mengtao Tang, Michael W. Topping, Lily Whitler\\
\textbf{Published:} 2026-04-23T10:23:19+00:00\\
\textbf{PDF:} \href{https://arxiv.org/pdf/2604.21516v1}{https://arxiv.org/pdf/2604.21516v1}\\
\textbf{Summary:}
JWST has revealed a population of super-luminous early galaxies with a volume density in excess of most expectations. The spectra reveal diverse properties: while some reveal strong emission lines characteristic of galaxies in the midst of strong bursts, others show weak emission lines that could reflect old stellar populations, large escape fractions, or post-burst star formation histories. Through the JWST Cycle 4 large program SPURS, we have obtained ultra-deep (29 hr) rest-frame UV spectroscopy of a z=9.3 super-luminous ($M_{\rm UV}=-21.66$) galaxy with large assembled stellar mass (1.6$\times$10$^9$ $M_\odot$) and extremely weak emission lines (H$\beta $ EW $\approx25$\textasciitilde{}Å). The strong stellar wind features and rest-optical line ratios suggest the galaxy is already significantly enriched, with a metallicity of 0.4--0.7\textasciitilde{}Z$_\odot$. The interstellar absorption lines reveal outflows ($v\simeq -161$\textasciitilde{}km\textasciitilde{}s$^{-1}$) with a large neutral gas covering fraction, suggesting that the weak emission lines are not due to large escape fractions. The combination of the Balmer break, weak emission lines, and stellar wind features constrains the star formation history, indicating a recent burst of star formation lasting 10--20 Myr followed by a downturn over the last 10\textasciitilde{}Myr. The observations suggest that $z\gtrsim 9$ weak emission line galaxies such as this source can be explained by stochastic star formation, provided that the downturns in star formation are recent (i.e., <10 Myr prior to observation). The ultra-deep grating spectrum enables the IGM damping wing to be characterized, decoupling the effects of local absorption. The smooth Ly$\alpha $ break indicates that this source, one of the most massive galaxies known at z>9, is likely situated in a small ionized bubble ($0.29_{-0.09}^{+0.11}$\textasciitilde{}pMpc), as is common at large neutral hydrogen fractions ($\bar{x}_{\rm HI}=0.81_{-0.21}^{+0.14}$).
\newpage
\subsection*{13. Signatures of Very Massive Stars in the Epoch of Reionization}
\textbf{Focus:} galaxies\\
\textbf{Authors:} Rui Marques-Chaves, Fabrice Martins, Daniel Schaerer, Miroslava Dessauges-Zavadsky, Ana Palacios\\
\textbf{Published:} 2026-04-23T09:54:38+00:00\\
\textbf{PDF:} \href{https://arxiv.org/pdf/2604.21493v1}{https://arxiv.org/pdf/2604.21493v1}\\
\textbf{Summary:}
We present ultra-deep ($\simeq 20-30$ hours), rest-frame UV spectroscopy with NIRSpec/JWST of two UV-bright galaxies at $z\sim 8.7$, CEERS-1019 and CEERS-1025 ($Z_{\rm neb} \simeq 0.1 Z_{\odot}$). The spectra reveal exceptionally strong P-Cygni profiles in wind lines (NV $\lambda $1240 and CIV $\lambda $1550) and significant broad and strong HeII $\lambda $1640 emission ($\rm EW\simeq 2-4$ A). We compare the observations with synthetic stellar population models at $Z_{\star} \simeq 0.1 Z_{\odot}$, both including and excluding very massive stars (VMS). Models including VMS provide a markedly improved fit to the data relative to non-VMS models ($\Delta $AIC and $\Delta $BIC $> 70$), which fail to reproduce the observed strengths of the wind features. A comparison with empirical spectra of VMS-dominated systems in the local Universe further supports this interpretation. The best-fit VMS models imply extremely young ages of the stellar populations ($\simeq 1.5-2.0$Myr) and high ionizing photon production efficiencies ($\log ξ_{\rm ion} [\rm Hz erg^{-1}] \gtrsim 25.8$), exceeding those inferred from non-VMS models by $\sim 0.1-0.2$ dex. These results provide evidence for an overabundance of VMS at high-$z$ with an IMF extending well beyond $100 M_{\odot}$, and highlight their potential role in shaping the rest-frame UV spectra, chemical enrichment, and ionizing output of galaxies in the early Universe.
\end{document}
